% Standalone section-ready fragment. Authoritative source:
% /workspace/shared/joint_realization_20261103/report.tex
% All claims below are for the joint layer-native implementation class there.

\section{An exact fixed-budget polynomial encoding}
\label{sec:budget-encoding}

Fix the finite normalized input data of the joint-native realization theorem.
In particular, the module set is nonempty; $J,E,S$ are finite and nonempty;
$A_i\one=p$ with $\sum_jp_j=1$; every $B_i$ and $Q$ is stochastic;
and the prescribed $w$ and $R$ have total mass one.
All input consistency and output-comparison squares are assumed checked.
The budget $K\in\mathbb Z_{\ge0}$ counts \emph{middle states}, not optimization
iterations. A middle-state budget $K$ is a total native-state budget $K+2$.

For each slot $v\in\{1,\ldots,K\}$ introduce real variables
\[
 d_{vh}\ (h\in E),\quad \eta_v,\qquad
 a_{jv},\rho_{jv},\sigma_{jv}\ (j\in J),\qquad
 b_{vs},\tau_{vs}\ (s\in S).
\]
There are $N=K(|E|+1+3|J|+2|S|)$ real variables. Write
$e_v=\sum_{h\in E}d_{vh}$ \emph{as an abbreviation}, not as an additional
unconstrained coordinate. The complete residual family is
\begin{align}
 D_{vh}&=d_{vh}^2-d_{vh}, &&(v,h),\label{eq:poly-selector}\\
 H_v&=1-e_v-\eta_v^2, &&(v),\label{eq:poly-cap}\\
 P_{jv}&=a_{jv}-\rho_{jv}^2, &&(j,v),\label{eq:poly-root-positive}\\
 G_{jv}&=e_v-a_{jv}-\sigma_{jv}^2, &&(j,v),\label{eq:poly-root-gate}\\
 T_{vs}&=b_{vs}-\tau_{vs}^2, &&(v,s),\label{eq:poly-row-positive}\\
 N_v&=\sum_sb_{vs}-e_v, &&(v),\label{eq:poly-row-normalization}\\
 C_{vix}&=\sum_{s:\pi_i(s)=x}b_{vs}
             -\sum_hd_{vh}B_i(h_i,x), &&(v,i,x\in S_i),\label{eq:poly-row-marginal}\\
 L_{iju}&=\sum_va_{jv}\sum_{h:h_i=u}d_{vh}-A_i(j,u),
             &&(i,j,u\in H_i),\label{eq:poly-local-root}\\
 W_s&=\sum_{j,v}a_{jv}b_{vs}-w_s, &&(s),\label{eq:poly-retained}\\
 O_{ju}&=\sum_{v,s}a_{jv}b_{vs}Q_{su}-R_{ju}, &&(j,u\in I).
             \label{eq:poly-observed}
\end{align}
The objective is the unweighted sum of squares of \emph{every scalar}
residual in this list, each included once:
\begin{equation}\label{eq:poly-objective}
 \Phi_K(x)=\sum_{v,h}D_{vh}^2+\sum_vH_v^2
 +\sum_{j,v}(P_{jv}^2+G_{jv}^2)+\sum_{v,s}T_{vs}^2
 +\sum_vN_v^2+\sum_{v,i,x}C_{vix}^2
 +\sum_{i,j,u}L_{iju}^2+\sum_sW_s^2+\sum_{j,u}O_{ju}^2.
\end{equation}
Every residual has degree at most two, so $\Phi_K$ is a nonnegative real
polynomial of degree at most four, and $-\nabla\Phi_K$ is polynomial of degree
at most three. All these are unconstrained real coordinates: inequalities
have been encoded by the square-slack equations, not silently omitted.

\begin{theorem}[Exact bounded-state realization]
\label{thm:budget-zero}
For $K\ge1$, the following are equivalent:
\begin{enumerate}
\item There is a joint layer-native implementation with at most $K$ middle states.
\item There exists $x\in\mathbb R^N$ with $\Phi_K(x)=0$.
\item $\min_{x\in\mathbb R^N}\Phi_K(x)=0$.
\item $\inf_{x\in\mathbb R^N}\Phi_K(x)=0$.
\end{enumerate}
In fact $\Phi_K$ is coercive, so its minimum is attained on every input.
If no such implementation exists, its minimum is strictly positive.
For $K=0$ there is no implementation and the empty-coordinate residual is
the strictly positive constant
$\sum_{i,j,u}A_i(j,u)^2+\sum_sw_s^2+\sum_{j,u}R_{ju}^2$.
\end{theorem}

\begin{proof}[Zero-set equivalence]
At a zero, \eqref{eq:poly-selector} gives $d_{vh}\in\{0,1\}$ and
\eqref{eq:poly-cap} gives $e_v\le1$. Consequently each slot has either no
selected tuple ($e_v=0$) or exactly one selected tuple $h(v)$ ($e_v=1$).
Equations \eqref{eq:poly-root-positive}--\eqref{eq:poly-root-gate} give
$0\le a_{jv}\le e_v$. Equations \eqref{eq:poly-row-positive} and
\eqref{eq:poly-row-normalization} give $b_{vs}\ge0$ and
$\sum_sb_{vs}=e_v$. Thus all $a_{jv}$ and $b_{vs}$ are zero in an inactive
slot. An active slot has a stochastic continuation row, and
\eqref{eq:poly-row-marginal} says
$B(v,-):=b_{v,-}\in\mathcal K_{h(v)}$.

Let $V$ be the active slots, put $A(j,v)=a_{jv}$, and set $f(v)=h(v)$.
Equations \eqref{eq:poly-local-root}, \eqref{eq:poly-retained}, and
\eqref{eq:poly-observed} are precisely the local root and global retained
and observed laws. The resulting $V$ is nonempty: for any one module,
summing \eqref{eq:poly-local-root} over $j,u$ gives
$\sum_{j,v}a_{jv}e_v=1$. Inactive slots vanish and active slots have $e_v=1$,
so $\sum_{j,v\in V}a_{jv}=1$. We have constructed a native implementation
with $|V|\le K$. Active slots with $\sum_ja_{jv}=0$ may also be deleted;
their presence is not needed to count the budget.

Conversely embed the states of any implementation with at most $K$ middle
states in the slots. For its occupied slots set $d_{v,f(v)}=1$, all other
selectors zero, $a_{jv}=A(j,v)$, and $b_{vs}=B(v,s)$. Its total root mass is
one, so $0\le a_{jv}\le1$. For unoccupied slots set all $d,a,b$ to zero.
Choose real square roots
\[
 \eta_v=\sqrt{1-e_v},\quad
 \rho_{jv}=\sqrt{a_{jv}},\quad
 \sigma_{jv}=\sqrt{e_v-a_{jv}},\quad
 \tau_{vs}=\sqrt{b_{vs}}.
\]
Every residual vanishes. The existence of these real square roots, rather
than their rationality, is all that is used. Coercivity and attainment are
proved below and establish the remaining equivalences.
\end{proof}

\subsection{Why empty coupling polytopes and unused slots cause no exception}
An active slot labelled $h$ supplies a row of $\mathcal K_h$, so an empty
$\mathcal K_h$ cannot be selected at a zero. An inactive slot supplies the
zero row, not a row normalized to one. It is therefore harmless even if
\emph{every} $\mathcal K_h$ is empty. In that case the target mass equations
make a zero impossible, as they should. A named unused local occurrence
need not have a preimage: no surjectivity has been added.

This also handles zero announcement mass, zero retained-state mass, and
zero tensor blocks. For example, if $p_j=0$, the nonnegative root equations
force $a_{jv}=0$ on every active slot; inactive slots already have zero root
weight. If $w_s=0$, then $a_{jv}b_{vs}=0$ for every $j,v$.
No division by one of these masses occurs in the encoding.

One must not impose a normalized row in $\mathcal K_h$ for every padded
factor and call the result a rank-at-most encoding. If $\mathcal K_h$ is
empty and the desired block is zero, its constrained rank is zero but that
padded normalization would be infeasible. The activity gate is what removes
this false obstruction.

\subsection{Exact relation to the constrained-rank resource law}
At a zero define
\[
 Z_{jhs}=\sum_{v:d_{vh}=1}a_{jv}b_{vs}.
\]
It is a feasible tensor for the preceding joint-native theorem. If $r_h$
is the number of active slots with label $h$, then
$Z^h$ factors through $r_h$ nonnegative columns and continuation rows in
$\mathcal K_h$. Thus
\[
 \sum_h\operatorname{rank}_{\mathcal K_h}(Z^h)\le\sum_hr_h\le K.
\]
Conversely, any feasible tensor with this constrained-rank sum at most $K$
provides factors, hence an implementation and a residual zero. Therefore
\begin{equation}\label{eq:budget-rank}
 \exists x:\Phi_K(x)=0
 \quad\Longleftrightarrow\quad
 \exists Z\in\mathcal P(D):
 \sum_h\operatorname{rank}_{\mathcal K_h}(Z^h)\le K.
\end{equation}
Here $\operatorname{rank}_{\mathcal K_h}(0)=0$ even for an empty polytope,
and a nonzero unfactorable block has rank $+\infty$.
Equivalently one can take the finite union over integer allocations
$r_h\ge0$ with $\sum_hr_h\le K$. The global-slot selectors encode that union
inside one polynomial system. They express an \emph{upper bound}; neither
nonzero incoming mass in every slot nor exact rank $K$ is required.
In particular $\operatorname{rank}(R)>K$ is an immediate obstruction, with
a nonzero $(K+1)$-minor providing an exact lower-bound witness when available.
The identity \eqref{eq:budget-rank} is about coupling-constrained
nonnegative factor rank, not ordinary matrix rank alone.

\subsection{Coercivity of the actual objective, without additional constraints}
\label{sec:budget-coercivity}
Fix $M\ge0$ and put $\varepsilon=\sqrt M$. On the sublevel set
$\Phi_K\le M$, the absolute value of every residual is at most
$\varepsilon$. Set
\[
 D_\varepsilon=\frac{1+\sqrt{1+4\varepsilon}}2.
\]
The selector residual gives $d_{vh}^2-d_{vh}\le\varepsilon$, hence
$|d_{vh}|\le D_\varepsilon$ and $e_v\ge-|E|D_\varepsilon$.
The cap residual gives
$e_v=1-\eta_v^2-H_v\le1+\varepsilon$ and therefore
\[
 \eta_v^2=1-e_v-H_v\le1+|E|D_\varepsilon+\varepsilon.
\]
The root positivity and gate residuals imply
\[
 -\varepsilon\le a_{jv}\le e_v+\varepsilon\le1+2\varepsilon,
 \qquad
 \rho_{jv}^2\le1+3\varepsilon,
 \qquad
 \sigma_{jv}^2\le1+3\varepsilon.
\]
Indeed $a=\rho^2+P\ge-\varepsilon$ and
$a=e-\sigma^2-G\le e+\varepsilon$; then
$\rho^2=a-P$ and $\sigma^2=e-a-G$ give the two bounds.
Similarly $b_{vs}\ge-\varepsilon$, while
$\sum_sb_{vs}=e_v+N_v\le1+2\varepsilon$. Writing $n=|S|$ and using the
lower bound for the other $n-1$ entries yields
\[
 -\varepsilon\le b_{vs}\le1+(n+1)\varepsilon,
 \qquad
 \tau_{vs}^2\le1+(n+2)\varepsilon.
\]
These estimates bound \emph{every} coordinate of the sublevel set; the
fixed row and target coefficients play no role in the bounds. Sublevel sets
are also closed, since $\Phi_K$ is continuous. They are therefore compact.
If a sequence with Euclidean norm tending to infinity had bounded objective,
it would lie in one such bounded sublevel set, a contradiction. Thus
$\Phi_K$ is coercive. Its nonempty sublevel at any evaluated point contains
a global minimizer by continuity and compactness. As $\Phi_K\ge0$,
infeasibility makes this attained minimum strictly positive.

This proof applies to the displayed gated encoding. A generic assertion
that every polynomial sum of squares attains its infimum would be false.
The gates and normalization equations are essential to this proof; dropping
them creates a different objective requiring a new argument.

\subsection{An optional redundant circle completion}
Every exact zero has $|x_q|\le1$ for every base coordinate $x_q$, as follows
from the exact selector, gate, and normalization equations. Consequently
one may introduce $z_q$ and add the residuals
\[
 U_q=x_q^2+z_q^2-1,\qquad
 \widetilde\Phi_K(x,z)=\Phi_K(x)+\sum_{q=1}^NU_q^2.
\]
This preserves zero feasibility, keeps degree at most four, and gives the
particularly short coercivity bound
\[
 \widetilde\Phi_K(x,z)
 \ge\frac{(\|(x,z)\|_2^2-N)^2}{N}\qquad(N>0).
\]
It follows from Cauchy--Schwarz applied to the $N$ numbers
$x_q^2+z_q^2-1$. It is not needed for the preceding coercivity theorem and
should not be represented as a missing assumption of it.

\subsection{The explicit positive-residual stationary point}
\label{sec:budget-stationary}
For any $K\ge1$, put
\[
 d=a=b=\rho=\sigma=\tau=0,\qquad \eta_v=1\quad(v=1,\ldots,K).
\]
All selector, cap, positivity, gate, normalization, and row-marginal
residuals vanish. The remaining residuals are $-A_i(j,u)$, $-w_s$, and
$-R_{ju}$. Their variable-dependent parts are bilinear in $(a,d)$ or
$(a,b)$, so their entire gradients vanish at this point. Hence
\[
 \nabla\Phi_K=0,\qquad
 \Phi_K=\sum_{i,j,u}A_i(j,u)^2+\sum_sw_s^2+\sum_{j,u}R_{ju}^2>0.
\]
This works even for an input that admits a zero. The fact that $\eta_v=1$
is nonzero causes no hidden gradient contribution: its residual $H_v$ is
zero, so the term $2H_v\nabla H_v$ is exactly zero.
For the optional circle completion set $z_q=1$ when $x_q=0$ and $z_q=0$
when $x_q=\eta_v=1$. Then every $U_q=0$, so all their gradient contributions
also vanish. Thus the completed objective has the same positive-residual
stationary obstruction.

More generally $\nabla\Phi_K=2J_F^{\mathsf T}F$, where $F$ is the residual
vector. The equation $J_F^{\mathsf T}F=0$ does not imply $F=0$. Along its
negative-gradient flow,
$d\Phi_K/dt=-\|\nabla\Phi_K\|^2$, but descent alone does not find a zero.
Coercivity bounds each trajectory in its initial sublevel and the smooth
polynomial vector field is locally Lipschitz, so forward solutions exist
for all time. No assertion that their limits have zero objective follows.
An optimization failure or a positive-residual stationary point is not an
infeasibility certificate.

\subsection{Inequalities, arithmetic, and exact certificates}
A weak inequality $g(x)\ge0$ is equivalent over the reals to the existence
of a slack $t$ with $g(x)-t^2=0$. Penalizing $g(x)^2$ instead imposes
$g(x)=0$, which is not the same condition. A strict inequality $g(x)>0$
can be encoded by $g(x)t^2-1=0$, but is unnecessary here and may destroy
compactness; exact positive mass in every padded slot should not be
smuggled into an at-most-budget question. If further polynomial weak
inequalities in the already bounded base coordinates are added using
square slacks, their slacks are bounded on sublevels too. Any more general
extension must recheck the degree and attainment arguments.

For rational input the displayed residual coefficients are rational.
Clearing denominators gives an existential sentence over the reals with
integer polynomial equations. Its size is polynomial in the expanded
finite input tables and $K$. If $K$ is binary encoded, it can first be
replaced by $\min(K,|J||E|)$, because the native reconstruction theorem
already bounds the needed number of middle states by $|J||E|$ on every
feasible instance. This gives membership in the existential theory of the
reals relative to explicit tables. It is not a polynomial-time algorithm,
a hardness proof, or a claim about succinctly specified product carriers.

The real-algebraic zero-set alternative gives two kinds of exact evidence:
\begin{itemize}
\item A point at which every residual is exactly zero is a feasibility
certificate and reconstructs the actual implementation as above. With
rational coefficients a nonempty real zero set has a point over the field
of real algebraic numbers: that field and $\mathbb R$ are real closed, and
quantifier elimination transfers the rational existential sentence between
them. Rational factor coordinates are \emph{not} asserted. Even square
slacks of rational factors need not be rational.
\item If the polynomial residuals are $F_1,\ldots,F_m$ and have no common
real zero, the real Nullstellensatz supplies polynomials $q_i,r_\ell$ with
\begin{equation}\label{eq:budget-nullstellensatz}
 -1=\sum_iq_iF_i+\sum_\ell r_\ell^2.
\end{equation}
At a putative zero the right side is nonnegative, so an exact polynomial
identity of this form certifies infeasibility. Conversely existence of
such an identity excludes a zero. This is a finite certificate principle,
not a promised small degree, a numerically stable search method, or a
rational-certificate-size claim. For rational input one can obtain a
certificate with real algebraic coefficients: fix the finite degree and
number of squares of a real certificate, express coefficient matching as
an existential rational system, and apply the same real-closed-field
transfer. Exact algebraic arithmetic then verifies it.
\end{itemize}
If the original linear tensor polytope is empty, its Farkas separator is a
simpler certificate excluding \emph{all} budgets. If that polytope is
nonempty but $K$ is too small, a general linear separator need not capture
the rank obstruction; the polynomial certificate principle applies.
A known positive global lower bound also certifies failure of the budget,
but a floating-point local optimizer value or a nonzero numerical residual
does not. Conversely, approximate residual zero without a rigorous error
bound or exact recovery does not prove existence. The instance-dependent
positive gap exists in every infeasible instance by coercivity, but this
does not supply its value or a universal numerical tolerance.

\subsection{Compositional bookkeeping}
To decompose the construction into components, give each scalar residual
one owner. Selector, gate, positivity, and row constraints can belong to
their individual slots; the local-root, retained, and observed sum
constraints are global constraints and should each be included once.
Shared coordinates can be identified through coordinate projections.
Fixed constants such as $A_i,B_i,w,R,Q$ belong in polynomial coefficients;
no claim that a constant-shift map is linear is needed. Duplicate
constraints with positive weights preserve the zero set but change the
objective and gradient field. Thus a particular compositional optimizer
must specify ownership and weights; existence of this exact encoding does
not supply a canonical choice of optimizer or a convergence theorem.
